\chapter{Cameras and perspective projection}
\label{chap:lecture07}

A range sensor reports a distance, whereas an ordinary camera produces an
image. The image contains information about the appearance and direction of
objects, but does not directly assign a three-dimensional position to each
pixel. To use a camera in robotics, we need a model connecting spatial points
to their image locations and an understanding of what information that mapping
discards.

This chapter first describes images as arrays of measurements, then develops
the pinhole camera model. The forward problem predicts where a known point
appears in an image. The reverse problem reveals why an image location alone
does not determine depth.

\section{Images as arrays of measurements}
Light from a scene reaches a photosensitive surface. A digital image sensor
contains an array of small sensing sites that convert incident light into
electrical signals. Readout electronics and analog-to-digital conversion
produce numerical values. Exposure, sensor response, and processing influence
those values; they are not direct measurements of object distance.

\subsection{Grayscale images}
A grayscale image can be represented as an $H\times W$ matrix, where $H$
is the number of rows and $W$ the number of columns. An 8-bit representation
stores values from 0 to 255. With a conventional grayscale display, 0 is
black, 255 is white, and intermediate values are shades of gray.
Figure~\ref{fig:l07-pixels} illustrates this correspondence using a small
image. A saturated value of 255 cannot distinguish stronger inputs beyond
the represented range.

\begin{figure}[htbp]
\centering\begin{tikzpicture}[scale=.8]
\foreach \row/\values in {0/{70,150,230,150,70},1/{150,230,255,230,150},2/{70,150,230,150,70}} {
\foreach \v [count=\col from 0] in \values {
\node at (\col+.5,-\row-.5) {$\v$};
\pgfmathsetmacro{\shade}{\v/255*100}
\fill[white!\shade!black] (\col+7,-\row) rectangle +(1,-1);
}}
\draw[step=1] (0,0) grid (5,-3);
\draw[step=1] (7,0) grid (12,-3);
\draw[axis] (5.3,-1.5)--(6.7,-1.5);
\node[above] at (2.5,0) {Stored intensities};
\node[above] at (9.5,0) {Grayscale display};
\end{tikzpicture}

\caption{A $3\times5$ array and its grayscale visualization. Each displayed
square corresponds to one matrix entry; larger values appear brighter.}
\label{fig:l07-pixels}
\end{figure}

Image dimensions are often written as width by height. Thus a
$640\times480$ image has 480 rows, 640 columns, and 307,200 pixels.
This convention reverses the order used to state matrix dimensions.
Keeping row, column, width, and height distinct prevents indexing errors.

\subsection{Color images and RGB values}
An RGB image assigns three values to each output pixel: red, green, and blue.
For 8-bit channels, some examples are
\begin{center}
\begin{tabular}{ll}
\toprule
RGB values & Displayed color\\
\midrule
$(255,0,0)$ & Red\\
$(0,255,0)$ & Green\\
$(0,0,255)$ & Blue\\
$(255,255,0)$ & Yellow\\
$(0,0,0)$ & Black\\
$(255,255,255)$ & White\\
\bottomrule
\end{tabular}
\end{center}
These are additive light mixtures. Mixing emitted red and green light is
different from mixing red and green pigments. A pixel such as
$(126,8,180)$ has strong blue and red contributions and appears purple.
Channel order is a convention: an application using BGR requires a different
interpretation of the same three stored numbers.

\subsection{Color filters and reconstructed pixels}
One common sensor design places a color-filter mosaic over the sensing
sites. A Bayer pattern repeats a $2\times2$ arrangement containing one red,
two green, and one blue filter. Figure~\ref{fig:l07-bayer} shows one possible
arrangement. Each site measures only its filtered channel; it does not
independently measure all three RGB components.

\begin{figure}[htbp]
\centering\begin{tikzpicture}[scale=.7]
\foreach \row in {0,1,2,3} {
\foreach \col in {0,1,2,3} {
\pgfmathtruncatemacro{\kind}{mod(\row,2)*2+mod(\col,2)}
\ifcase\kind\def\shade{red!40}\def\letter{R}
\or\def\shade{green!35}\def\letter{G}
\or\def\shade{green!35}\def\letter{G}
\or\def\shade{blue!35}\def\letter{B}\fi
\filldraw[fill=\shade] (\col,-\row) rectangle +(1,-1);
\node at (\col+.5,-\row-.5) {\letter};
}}
\node[above] at (2,0) {One channel sampled per site};
\draw[axis] (4.4,-2)--(6.5,-2) node[midway,above] {demosaic};
\node[draw,align=center,minimum height=1.8cm] at (8.3,-2) {RGB image\\three values\\per output pixel};
\end{tikzpicture}

\caption{An RGGB Bayer mosaic. Each letter denotes the channel sampled at
one sensing site. Demosaicing estimates missing channels to form an RGB image.}
\label{fig:l07-bayer}
\end{figure}

The missing color components are estimated using neighboring samples in a
process called \emph{demosaicing}. A repeated four-site pattern should
therefore not be confused with four independent full-color measurements or
with a universal rule that four sites produce exactly one output pixel.
The raw mosaic and the reconstructed RGB image are different data
representations.\footnote{See
\href{https://docs.baslerweb.com/demosaicing-mode}{Basler, Demosaicing Mode}.}

\section{The pinhole camera model}
An ideal pinhole camera consists of a light-tight enclosure, a small aperture,
and an image plane behind it. In geometric optics, the ray from a scene point
that passes through the aperture intersects the image plane at the point's
image. Rays cross at the aperture, so the physical image is inverted.

The model treats the aperture as a single \emph{projection center}. Actual
cameras use finite apertures and lenses to collect and focus light. The
pinhole model describes their central projection approximately; it does not
model blur, diffraction, or lens distortion. Making a physical pinhole ever
smaller is not an unlimited route to a sharper image, because light throughput
and diffraction also matter.

Let $C$ be the camera frame with origin $O_C$ at the projection center.
Use the optical-axis convention
\[
 z_C:\text{ forward into the scene},\qquad
 x_C:\text{ image-right},\qquad y_C:\text{ image-down}.
\]
These directions form a right-handed frame. In a two-dimensional section,
the page may be rotated for readability without changing the frame definition.
Write a scene point as
\[
 B^C=\begin{bmatrix}X\\Y\\Z\end{bmatrix},\qquad Z>0.
\]
The variables $X,Y,Z$ are the point's coordinates in $C$, measured in a
common length unit. The quantity $Z$ is depth along the optical axis, not
generally the Euclidean distance $\|B^C\|$.

\subsection{Physical and virtual image planes}
Place the physical image plane at $z_C=-f$, with $f>0$. The line through
$O_C$ and $B^C$ has coordinates $\lambda[X\;Y\;Z]\trans$.
Its intersection with the physical plane has $\lambda=-f/Z$, giving
\[
 x_{\mathrm{physical}}=-f\frac{X}{Z},\qquad
 y_{\mathrm{physical}}=-f\frac{Y}{Z}.
\]
The negative signs express the inversion when image axes are parallel to
the corresponding camera axes.

For calculations, introduce a \emph{virtual image plane} at $z_C=f$,
in front of the projection center. Its coordinates are
\begin{equation}
 \boxed{x'=f\frac{X}{Z},\qquad y'=f\frac{Y}{Z}.}
 \label{eq:l07-metric}
\end{equation}
This is a modeling convention, not a physical sensor placed in front of the
camera. It gives an upright image and removes the two negative signs.

\begin{figure}[htbp]
\centering\begin{tikzpicture}[scale=.95]
\draw[axis] (-2.7,0)--(6.5,0) node[right] {$z_C$};
\draw[axis] (0,0)--(0,-2.4) node[below] {$x_C$};
\draw[very thick] (-2,-1.9)--(-2,1.9);
\draw[thick,sensorframe] (2,-1.9)--(2,1.9);
\node[above,align=center] at (-2,1.9) {Physical plane\\$z_C=-f$};
\node[above,align=center,sensorframe] at (2,1.9) {Virtual plane\\$z_C=f$};
\draw[thick,robotframe] (-2,.6)--(5,-1.5);
\fill (0,0) circle(2pt) node[above right] {$O_C$};
\fill[robotframe] (5,-1.5) circle(2pt) node[right] {$B=(X,Y,Z)$};
\fill[robotframe] (-2,.6) circle(2pt);
\fill[sensorframe] (2,-.6) circle(2pt) node[right] {$x'=fX/Z$};
\draw[projection] (5,-1.5)--(5,0);
\draw[<->] (0,-2)--(2,-2) node[midway,below] {$f$};
\draw[<->] (-2,-2)--(0,-2) node[midway,below] {$f$};
\node[left] at (-2,.6) {$-fX/Z$};
\end{tikzpicture}

\caption{A section through the $x_Cz_C$ plane, drawn with positive $x_C$
down the page. The physical and virtual image coordinates have opposite
signs. The $y_C$ coordinate is omitted from this section.}
\label{fig:l07-pinhole}
\end{figure}

The same relationship follows from similar triangles. In this section,
$X=\ell\sin\theta$ and $Z=\ell\cos\theta$, so $X/Z=\tan\theta$.
For the virtual image triangle, $x'/f=\tan\theta$. Equating the ratios
gives $x'=fX/Z$. Repeating the construction in the $y_Cz_C$ plane gives
the second equation.

The parameter $f$ is the projection distance in this ideal model. For a
lens camera, an effective calibrated model is used; $f$ should not simply
be identified with the distance from the front of the camera housing to
the sensor. Changing focus or zoom can change the appropriate calibration.

\section{From metric image coordinates to pixels}
The coordinates $x',y'$ are lengths measured from the optical-axis
intersection with the virtual plane. Image processing instead uses pixel
coordinates measured relative to an image origin. Two operations are needed:
convert lengths to pixel units and shift the origin.

Let $\alpha_x,\alpha_y$ be the pixel densities along the two image axes,
with units pixels per meter when $f$ is in meters. Define
\[
 f_x=\alpha_x f,\qquad f_y=\alpha_y f,
\]
which are focal lengths expressed in pixel units. Square sampling gives
equal densities; the general model permits different scales.
Let $(c_x,c_y)$ be the \emph{principal point}, the image coordinates of the
optical-axis intersection. The projection equations are
\begin{equation}
 \boxed{u=f_x\frac{X}{Z}+c_x,\qquad
 v=f_y\frac{Y}{Z}+c_y.}
 \label{eq:l07-projection}
\end{equation}
Here $u$ increases to the right and $v$ downward. They correspond to the
horizontal and vertical image coordinates, also denoted $x_B^I$ and $y_B^I$.
The formulas assume perpendicular image axes and no distortion.

\subsection{Image center and indexing conventions}
For an ideal centered camera, a convenient continuous coordinate convention
places the image boundaries at $u=0,W$ and $v=0,H$. Then
\[
 c_x=\frac W2,\qquad c_y=\frac H2.
\]
Under this convention the pixel centers are at half-integer coordinates.
If instead zero-based integer coordinates label pixel centers, the geometric
center of the array is $((W-1)/2,(H-1)/2)$. These conventions differ by half
a pixel; either can be used consistently. In a calibrated camera, the
principal point need not equal the exact geometric center.

A projection generally gives continuous coordinates, not an integer array
index. Selecting a pixel or interpolating an image requires a sampling rule.
In row--column array notation, the vertical coordinate selects the row and
the horizontal coordinate the column. A pixel location is therefore addressed
in the order $[\text{row},\text{column}]$, corresponding to $[v,u]$ after
the chosen conversion to indices.

\subsection{What calibration provides}
Camera calibration estimates the model parameters from observations of known
geometry, such as a checkerboard photographed in several poses. In practice,
calibration commonly estimates $f_x,f_y,c_x,c_y$ directly and may also
estimate lens-distortion parameters. The separate physical quantities $f$
and $\alpha_x,\alpha_y$ are not needed to use
Equation~\eqref{eq:l07-projection}.\footnote{See
\href{https://docs.opencv.org/doc/doxygen/html/d4/d93/group__calib.html}{OpenCV, Camera Calibration}.}

The pinhole equations apply to an ideal or appropriately undistorted image.
Using them on raw images with substantial unmodeled distortion introduces
systematic position errors. A detailed calibration procedure is outside
the scope of this chapter.

\section{Worked projection example}
Suppose a camera has
\[
 W=640,\quad H=480,\quad
 \alpha_x=\alpha_y=5\times10^4\,\mathrm{px/m},\quad
 f=3\times10^{-3}\,\mathrm m.
\]
Use the centered continuous-image convention, so $(c_x,c_y)=(320,240)$.
The focal lengths in pixels are
\[
 f_x=f_y=(5\times10^4)(3\times10^{-3})=150\,\mathrm{px}.
\]
For $B^C=[2\;0\;10]\trans\,\mathrm m$,
\begin{align}
 u&=(150\,\mathrm{px})\frac{2\,\mathrm m}{10\,\mathrm m}
       +320\,\mathrm{px}=350\,\mathrm{px},\\
 v&=(150\,\mathrm{px})\frac{0\,\mathrm m}{10\,\mathrm m}
       +240\,\mathrm{px}=240\,\mathrm{px}.
 \label{eq:l07-example}
\end{align}
The ratios are dimensionless, leaving pixel units in both equations. The
point appears 30 pixels to the right of the principal point and at the same
vertical coordinate, as shown in Figure~\ref{fig:l07-image}.

\begin{figure}[htbp]
\centering\begin{tikzpicture}[x=.012cm,y=-.012cm]
\draw[thick] (0,0) rectangle (640,480);
\draw[axis] (0,0)--(690,0) node[right] {$u$};
\draw[axis] (0,0)--(0,520) node[below] {$v$};
\draw[projection] (320,0)--(320,480);
\draw[projection] (0,240)--(640,240);
\fill (320,240) circle(2pt) node[below left] {$(320,240)$};
\fill[robotframe] (350,240) circle(2pt) node[above right] {$B:\ (350,240)$};
\node[above] at (320,-10) {$W=640$};
\node[right] at (650,380) {$H=480$};
\node[above left] at (0,0) {$(0,0)$};
\end{tikzpicture}

\caption{The projected location $(350,240)$ in the $640\times480$ example.
Dashed lines pass through the principal point. Coordinates refer to the
continuous image plane, before selecting discrete pixel indices.}
\label{fig:l07-image}
\end{figure}

\subsection{When does a projected point appear in an image?}
The model assumes $Z>0$. At $Z=0$ the division is undefined, while a point
with $Z<0$ is behind a forward-looking camera even if its algebraic projection
lies within the image bounds. A point in front must also project inside the
image area. In the continuous boundary convention, this means
$0\le u<W$ and $0\le v<H$ for sampling within the image.

These checks establish geometric eligibility, not guaranteed visibility:
another surface can occlude the point. Detecting the desired object in an
image is also a separate task from computing where a known point projects.

\section{Why one image point does not determine depth}
Rearranging the projection equations gives
\begin{equation}
 \frac XZ=\frac{u-c_x}{f_x},\qquad
 \frac YZ=\frac{v-c_y}{f_y}.
 \label{eq:l07-inverse}
\end{equation}
The right sides are known for a calibrated camera and measured image point,
but the left sides are ratios. Two equations constrain three spatial
coordinates. Writing the result as a family of points makes the remaining
ambiguity explicit:
\begin{equation}
 \boxed{B^C=Z\begin{bmatrix}(u-c_x)/f_x\\(v-c_y)/f_y\\1\end{bmatrix},
 \qquad Z>0.}
 \label{eq:l07-ray}
\end{equation}
The image determines a ray from the projection center. It does not specify
where along that ray the object lies. The vector in brackets is not generally
a unit vector: its scale factor here is axial depth $Z$, not range.

\begin{figure}[htbp]
\centering\begin{tikzpicture}
\draw[axis] (0,0)--(7,0) node[right] {$z_C$};
\draw[thick,sensorframe] (1.5,-.6)--(1.5,2.3);
\node[above,sensorframe] at (1.5,2.3) {Virtual image plane};
\draw[axis,robotframe] (0,0)--(6.5,2.6);
\fill (0,0) circle(2pt) node[below] {$O_C$};
\fill[sensorframe] (1.5,.6) circle(2pt) node[below right] {one image point};
\fill[robotframe] (3,1.2) circle(2pt) node[above] {$B_1$};
\fill[robotframe] (5,2) circle(2pt) node[above] {$B_2$};
\end{tikzpicture}

\caption{Distinct points on one viewing ray have the same image coordinates.
Projection removes their common spatial scale.}
\label{fig:l07-ray}
\end{figure}

In the worked example, $(u,v)=(350,240)$ gives $X/Z=0.2$ and $Y/Z=0$.
The points $[1\;0\;5]\trans$, $[2\;0\;10]\trans$, and
$[4\;0\;20]\trans$ all produce that same location. More generally,
multiplying $X,Y,Z$ by any positive constant leaves their ratios unchanged.
This geometric argument explains the ambiguity beyond simply counting
equations and unknowns.

\section{Additional information for recovering position}
Recovering a three-dimensional point requires information beyond one
unconstrained image location. Several sources are possible.

\subsection{Two cameras}
A stereo pair observes the same scene from distinct positions. If a point
is matched between the images and the cameras' calibration and relative
pose are known, the two viewing rays can be used to triangulate its position.
The additional image coordinates are useful because they come from another
viewpoint, not merely because more equations have been written down.
Degenerate geometry and measurement noise can prevent a unique or accurate
intersection. Stereo also requires identifying the same scene point in
both images.

\subsection{Known object geometry}
Known size can supply scale information. For example, a segment of known
length $L$ parallel to the image plane at constant depth $Z$ spans
\[
 s=f_x\frac LZ
\]
pixels when aligned with the horizontal axis. Hence $Z=f_xL/s$.
This simple relation assumes the stated alignment and a known correspondence
between the measured image endpoints and the object's endpoints. Unknown
orientation or occlusion makes an apparent-size argument more complicated.

\subsection{A known surface or depth}
If $Z$ is known independently, Equation~\eqref{eq:l07-ray} determines the
point directly. A downward-looking camera above a known level ground plane
provides one useful special case. When its optical axis is perpendicular
to the ground, points on that plane have axial depth equal to the camera's
height above it. If the camera is tilted, height is not generally equal to
each point's $Z$; the viewing ray must instead be intersected with the
ground plane using the camera pose. A point above the ground does not satisfy
the ground-plane assumption.

These methods add hardware, known scene geometry, or constraints on the
task. They do not remove the need to state which extra information makes
depth recovery possible.

\section{Connecting camera geometry to robot frames}
The camera optical frame usually differs from a robot frame whose forward
axis is $x_R$. The mounting transform includes both the position of the
camera origin and the orientation of its axes. Once a point's camera-frame
position has been obtained, its world coordinates are
\begin{equation}
 B^W=R_C^W B^C+p_C^W,\qquad
 H_C^W=H_R^W H_C^R.
 \label{eq:l07-world}
\end{equation}
The mounting calibration $H_C^R$ is distinct from the internal projection
parameters $f_x,f_y,c_x,c_y$.

Conversely, to predict the image of a known world point, first compute
\[
 B^C=(R_C^W)\trans(B^W-p_C^W),
\]
then apply Equation~\eqref{eq:l07-projection}. Coordinate transformation
and perspective projection are different operations: the former preserves
the spatial point, while the latter divides by depth and discards scale.
Transforming a viewing ray into the world frame does not by itself resolve
the missing depth.

\clearpage
\section{Additional exercises}
Use the virtual image plane, zero distortion, and the stated continuous-image
convention. Keep depth and Euclidean range distinct.
\begin{enumerate}
\item \textbf{Image representation.} How many rows, columns, and pixels are
in a $640\times480$ image? How many bytes store an uncompressed 8-bit grayscale
image and an 8-bit-per-channel RGB image, ignoring padding and metadata?

\item \textbf{Color sampling.} Sketch a $4\times4$ RGGB Bayer mosaic.
How many sites measure each channel? Explain why each raw site does not
already contain an RGB triple and what demosaicing supplies.

\item \textbf{Projection and visibility.} With $f_x=f_y=150$ and principal
point $(320,240)$, project $[2\;-1\;5]\trans$ and $[12\;0\;5]\trans$.
Which lies inside a $640\times480$ image? Explain why the algebraic projection
of $[0\;0\;-5]\trans$ does not establish visibility.

\item \textbf{Units and scale.} Let $f=4\,\mathrm{mm}$ and
$\alpha_x=\alpha_y=10^5\,\mathrm{px/m}$. Find $f_x,f_y$.
Using principal point $(320,240)$, project $[1\;0.5\;4]\trans\,\mathrm m$.
What changes if all three point coordinates are doubled?

\item \textbf{Back-projection.} For $f_x=f_y=150$, principal point $(320,240)$,
and image location $(350,255)$, write the full viewing ray.
Find its point at depth $Z=8\,\mathrm m$ and compute that point's range.
Explain why the image alone cannot select this depth.

\item \textbf{Known size.} A horizontal segment of length $0.4\,\mathrm m$
lies parallel to the image plane and spans 20 pixels with $f_x=150$.
Find its depth. Name two assumptions that could fail when this method is
applied to an arbitrary object in a photograph.

\item \textbf{Camera and world frames.} A camera has $R_C^W=I_3$ and
$p_C^W=[1\;0\;2]\trans\,\mathrm m$. For the world point
$B^W=[3\;0\;12]\trans\,\mathrm m$, compute $B^C$ and project it using
$f_x=f_y=150$, $(c_x,c_y)=(320,240)$. Explain what additional information
would be needed to reverse this calculation from the image point alone.
\end{enumerate}
